
\documentclass[11pt]{article}
\usepackage[margin=1in]{geometry}
\usepackage{amsmath,amssymb,booktabs,hyperref,graphicx,parskip,enumitem}
\usepackage{xcolor}
\hypersetup{colorlinks=true,linkcolor=blue!50!black,urlcolor=blue!50!black}
\setlist{nosep}

\title{Replication Report: Exactly-Solved Model of Light-Scattering Errors in Quantum Simulations with Metastable Trapped-Ion Qubits\\[2pt]\large OSTI 2441075 \textbar{} Coverage 8/10, Agreement 9/10}
\author{Rick Stevens \and Ollie (AI Assistant)\\\small Argonne National Laboratory}
\date{2026-04-27}

\begin{document}\maketitle

\section{Source paper}
\begin{itemize}
\item \textbf{Title:} Exactly-Solved Model of Light-Scattering Errors in Quantum Simulations with Metastable Trapped-Ion Qubits
\item \textbf{Authors:} Foss-Feig, Vikas, et al.
\item \textbf{Year:} 2024
\item \textbf{Domain:} Quantum Information / AMO
\item \textbf{Identifier:} OSTI 2441075
\end{itemize}


\section{Replication objective}
The central claim of the paper concerns quantum information / amo. Our objective was to reproduce the headline numerical/qualitative results within the constraints of the AI-driven Replication Project (24--72\,h, commodity GPU/CPU compute, no author contact). A replication plan is archived at \texttt{2441075-Exactly-solved-model-of-light-scattering-errors-in/replication\_plan.pdf}.

\section{Methodology}
We followed the standard Replication Project workflow: ingest the source PDF, extract method and data pointers, draft a replication plan, attempt the smallest end-to-end pipeline that produces a comparable number, then scale up where compute and data permit. Tools used in this replication are catalogued in the meta-paper's computational tools inventory (see \texttt{COMPUTATIONAL\_TOOLS\_INVENTORY.md}).



\subsection*{Paper-directory README excerpt}
\small
\par\medskip\textbf{Exactly-solved model of light-scattering errors in quantum simulations with metastable trapped-ion qubits}\par
\par$\bullet$\ **OSTI ID:** 2441075
\par$\bullet$\ **Rank:** \#24
\par$\bullet$\ **Replication Score:** 9/10
\par$\bullet$\ **Open-Source Tools:** Yes
\par$\bullet$\ **Code Repository:** No
\par$\bullet$\ **Tools:** Python (implied for numerical integration), Simpson's method (for ODEs)

\par\medskip\textbf{\# Why This Paper}\par
Analytic/numerical, all equations/algorithms specified, synthetic data, and quantitative validation. No code repo but trivial for AI.

\par\medskip\textbf{\# Replication Plan}\par
AI extracts formulas/algorithms, implements Lindbladian master equation, generates synthetic data, and reproduces results.

\par\medskip\textbf{\# Status}\par
\par$\bullet$\ [ ] Paper reviewed
\par$\bullet$\ [ ] Code/tools identified
\par$\bullet$\ [ ] Code implemented/cloned
\par$\bullet$\ [ ] Results reproduced
\par$\bullet$\ [ ] Results validated against paper
\normalsize

The replication was driven by an OpenClaw agent loop (Claude Opus / GPT-5 class models) issuing shell, Python, and (where applicable) GPU jobs. Compute usage and wall-time for individual papers are summarised in the meta-paper's resource table; not separately recorded for this paper unless noted in the Tier-Lift artefact. Sub-agents were spawned for replication-plan generation and (for papers in batches 1--4) for tier-lift extension runs.

\section{What was reproduced}
\begin{itemize}
\item Full analytic solution Eqs. 6-10 with all four I, R, L, B components (including the paper's novel L and B terms)
\item 5-jump-operator Lindblad numerical master equation with correct factor-of-2 jump convention
\item ⁴⁰Ca⁺ branching ratios (94.5\% leakage / 3.9\% Raman / 1.6\% elastic)
\item Selection rule: |1⟩=|mJ=-5/2⟩ dark under π polarization
\item Machine-precision analytic vs numerical agreement (10⁻¹¹–10⁻¹³)
\item Figure 2 GHZ-fidelity qualitative trends (N dependence, postselection gain, B-field effect)
\item Figure 3 correlation decay (faster than exp(-mΓt)) and perpendicular-correlation growth
\item Figure 4 spin-squeezing ξ² ∝ N\textasciicircum{}(-2/3) with scattering floor
\end{itemize}

\section{What was missing or incomplete}
\begin{itemize}
\item Penning-trap 2D triangular-lattice ion-crystal mode structure
\item Full 3D normal-mode J\_ij computation (we used 1D chain)
\item Laser power P and beam waist w₀ (not specified in paper)
\item Exact trap frequencies ωz, ωr (referenced to Ref [20])
\item GHZ fidelity fourth-order δJ\_ij correction terms
\item Quantitative N-thresholds from Fig 2 match
\item Experimental validation against real Ba⁺ or Ca⁺ data
\end{itemize}
The dominant blocker categories for this paper, mapped onto the meta-paper's 9-category friction taxonomy (\S4), are listed in \S\ref{sec:friction} below.

\section{Quantitative agreement}
Per-quantity numerical comparisons are not separately tabulated in the structured evaluation record for this paper beyond the agreement rationale below; where the rationale references specific numbers, they are reproduced verbatim. A full numerical comparison table is recorded in the per-replication artefacts (see \S\ref{sec:repro}). For papers below 8/8, the agreement rationale identifies the specific quantities that diverge.

\paragraph{Agreement rationale.} Analytic-numerical agreement at machine precision (errors 10⁻¹¹–10⁻¹³ across all tested cases). Trace preservation to 5.5×10⁻¹⁶, positivity to 10⁻¹⁸ — essentially exact. All qualitative figure features correct: fidelity ↓ with N, Postselection helps, m-body correlations decay faster than exp(-mΓt), ξ² ∝ N\textasciicircum{}(-2/3) scaling with scattering floor. Quantitative figure values differ because paper's Penning-trap J\_ij differs from our 1D-chain J\_ij.

\section{Score and friction-point tagging}\label{sec:friction}
\begin{itemize}
\item \textbf{Coverage:} 8/10 --- Complete implementation of analytic solution (Eqs. 6-10) with all four terms I, R, L, B — including the paper's novel leakage (L) and combined (B) contributions. Five-jump-operator Lindblad master equation implemented numerically with correct factor-of-2 jump convention. ⁴⁰Ca⁺ branching ratios (94.5\% leakage / 3.9\% Raman / 1.6\% elastic). Figures 1 (schematic), 2 (GHZ fidelity), 3 (correlations), 4 (squeezing) reproduced qualitatively. Missing: full Penning-trap 2D crystal mode computation (we used 1D chain approximation for J\_ij).
\item \textbf{Agreement:} 9/10 --- Analytic-numerical agreement at machine precision (errors 10⁻¹¹–10⁻¹³ across all tested cases). Trace preservation to 5.5×10⁻¹⁶, positivity to 10⁻¹⁸ — essentially exact. All qualitative figure features correct: fidelity ↓ with N, Postselection helps, m-body correlations decay faster than exp(-mΓt), ξ² ∝ N\textasciicircum{}(-2/3) scaling with scattering floor. Quantitative figure values differ because paper's Penning-trap J\_ij differs from our 1D-chain J\_ij.
\end{itemize}

\noindent\textbf{Friction points encountered (taxonomy tags):}
\begin{itemize}
\item \textbf{F7}: Missing surrogate calculators (xTB/DFT/closed solver)
\end{itemize}

\section{Follow-on questions}
\begin{itemize}
\item Q1: Implement the full 2D Penning-trap equilibrium-crystal mode solver and recompute J\_ij — do the Fig 2 fidelity crossovers match the paper's reported N\_threshold at B=0.9T and B=4.5T within 10\%?
\item Q2: Extend the exact solution from π to σ± polarization where |1⟩ is no longer dark — derive new analytic expressions and quantify the fidelity penalty of losing the dark-state protection.
\item Q3: Test the exact solution against a full Lindblad simulation at N=20 (near the boundary of tractable numerics) — does the product-form correlation factorization still agree to machine precision, or do many-body corrections emerge?
\item Q4: For ¹³⁸Ba⁺ (different branching ratios, \textasciitilde{}70\% leakage), do the novel L and B terms still dominate the error budget, and does postselection recover the same 5.5\% residual-error advantage?
\item Q5: Combine the exact scattering model with realistic single-qubit gate errors (miscalibration, Stark shifts) to produce a full end-to-end error budget for a 50-qubit GHZ preparation circuit — at what N does gate error overtake scattering as the dominant error source?
\end{itemize}

\section{Reproducibility pointers}\label{sec:repro}
\begin{itemize}
\item \textbf{Paper directory:} \texttt{2441075-Exactly-solved-model-of-light-scattering-errors-in}
\item \textbf{Replication artefacts:}
\begin{itemize}
\item \texttt{/Users/stevens/Dropbox/REPLICATE-PROJECT/2441075-Exactly-solved-model-of-light-scattering-errors-in/replication}
\end{itemize}
\item \textbf{Replication-dir hint from JSONL:} \texttt{\textasciitilde{}/Dropbox/REPLICATE-PROJECT/2441075-Exactly-solved-model-of-light-scattering-errors-in/replication/}
\item \textbf{Status:} Not recorded
\end{itemize}

To re-run, change directory into the paper directory above and consult its \texttt{README.md} for the entry-point script. Most replications follow the convention \texttt{replication/run\_*.sh} or \texttt{replication/<subdir>/run.py}.

\end{document}
